\def\stat{ch-sdn}





\def\tit{БАЛАНСИРОВКА НАГРУЗКИ В ЗАЩИЩЕННЫХ СЕТЯХ 


С~ИСПОЛЬЗОВАНИЕМ ТЕХНОЛОГИИ SDN}





\def\titkol{Балансировка нагрузки в~защищенных сетях 


с~использованием технологии SDN}








\def\aut{О.\,Ю.~Гузев$^1$, И.\,В.~Чижов$^2$}





\def\autkol{О.\,Ю.~Гузев, И.\,В.~Чижов}





\titel{\tit}{\aut}{\autkol}{\titkol}





\index{Гузев О.\,Ю.}


\index{Чижов И.\,В.}


\index{Guzev O.\,Yu.}


\index{Chizhov I.\,V.}





%{\renewcommand{\thefootnote}{\fnsymbol{footnote}}


%\footnotetext[1]


%{Работа выполнена при поддержке РФФИ (проект 15-07-02053).}}





\renewcommand{\thefootnote}{\arabic{footnote}}


\footnotetext[1]{Центр научных исследований и~перспективных разработок, ОАО <<ИнфоТеКС>>, 


\mbox{oleg.guzev@infotecs.ru}}


\footnotetext[2]{Московский государственный университет им.\ М.\,В.~Ломоносова, факультет вычислительной математики 


и~кибернетики; Институт проб\-лем информатики Федерального исследовательского центра <<Информатика 


и~управ\-ле\-ние>> Российской академии наук, \mbox{ichizhov@cs.msu.ru}}








 \label{st\stat}





                  \thispagestyle{headings}


                  


  %\vspace*{-8pt}











\Abst{Технология SDN 


(Software-Defined 


Networking, про\-грам\-мно-кон\-фи\-гу\-ри\-ру\-емые сети)


по сравнению с~традиционными IP-се\-тя\-ми позволяет 


программировать поведение сети при помощи централизованного контроллера. Передающие 


трафик устройства в~этом случае занимаются только пересылкой фреймов на основании 


таблиц потоков, загружаемых в~них контроллером. Таблица потоков строится на 


контроллере в~ходе обработки информации о~приходящих на передающие устройства 


потоках трафика. Перечисленные свойства технологии были использованы при создании 


SDN-ба\-лан\-си\-ров\-щи\-ка нагрузки для устройств защищенных сетей. В~статье 


рассматривается архитектура и~программное обеспечение балансировщика. Приводятся 


описания схем и~результаты экспериментов по балансировке нагрузки на такие устройства, 


как L3-крип\-то\-шлюз, TLS (Transport Layer Security, защита транспортного


уровня) шлюз, IDS


(Intrusion Detection System, система обнаружения вторжений).}





\KW{SDN; ПКС; программно-конфигурируемые сети; контроллер; OpenFlow; криптошлюз; 


TLS; система обнаружения вторжений; IDS; балансировка нагрузки; DPDK; Open vSwitch; 


Beacon}





\DOI{10.14357\!/\!08696527180110}








\vskip 10pt plus 9pt minus 6pt








      \thispagestyle{headings}


      


     % \vspace*{-10pt} 





\section{Введение}





  Традиционные IP-сети в~большинстве своем комплексны, статичны 


  и~сложны в~обслуживании~[1--4]. При этом с~каждым годом бизнес предъявляет все 


большие требования к~их функционалу и~производительности~\cite{5-sdn}. 


Такое противоречивое состояние развития IP-се\-тей требует применения 


новых подходов, поз\-во\-ля\-ющих упростить, стандартизировать, 


виртуализировать, централизовать сетевые устройства (СУ), процессы и~управление 


ими. Одним из таких подходов является технология SDN~\cite{2-sdn, 6-sdn}.


  


  В данной статье рассматривается использование технологии SDN для 


создания высокопроизводительного, конфигурируемого балансировщика 


нагрузки и~применение данного балансировщика для масштабирования таких 


устройств защищенных сетей, как L3-крип\-то\-шлюз, TLS-шлюз, IDS.


  


  Данная задача является актуальной, поскольку к~балансировщику нагрузки, 


используемому в~указанной сфере, помимо поддержки базовых алгоритмов 


балансировки и~производительности предъявляются дополнительные 


специфические требования с~точки зрения функциональности. К~таким 


требованиям можно отнести:


  \begin{itemize}


\item возможность изменения содержимого полей служебных заголовков 


Ethernet-фрей\-ма, IP-па\-ке\-та, TCP\!/\!UDP (Transmission


Control Protocol/User Datagram Protocol) сег\-мента;


\item поддержку протоколов туннелирования GRE (Generic Routing Encapsulation), VXLAN


(Virtual eXtensible Local Area Network);


\item согласованную работу с~балансировщиком нагрузки на другом конце 


масштабируемого защищенного канала связи для обслуживания случаев 


отказа отдельных линков;


\item использование любого набора полей служебных заголовков 


фрей\-ма\!/\!па\-ке\-та\!/\!сегмента для идентификации балансируемого потока трафика.


\end{itemize}





  Кроме того, с~точки зрения пути прохождения полезного трафика можно 


выделить несколько типов СУ, для которых может 


требоваться балансировка нагрузки (рис.~1).


  





\begin{figure}[b] %fig1


   \vspace*{1pt}


 \begin{center}


 \mbox{%


 \epsfxsize=126.384mm 


 \epsfbox{chsd-1.eps}


 }


\end{center}


 \vspace*{-11pt}


\Caption{Типы СУ, для которых выполняется балансировка нагрузки: 


(\textit{а})~СУ устанавливаются <<в~разрыв>> пути прохождения полезного 


трафика;  


(\textit{б})~балансировщик нагрузки <<заворачивает>> полезный трафик на 


СУ, после обработки на которых трафик возвращается на балансировщик для 


дальнейшей коммутации; (\textit{в})~полезный трафик <<зеркалируется>> 


балансировщиком на СУ для анализа}


\end{figure}


  


  Большинство балансировщиков нагрузки традиционных IP-се\-тей не 


отвечают всем указанным требованиям. Кроме того, в~них отсутствует 


возможность программируемости, что затрудняет использование одного 


устройства в~разных сценариях балансировки. В~то же время технология SDN 


и~протокол OpenFlow легко позволяют реализовать все перечисленные 


требования.


  


\section{Программно-аппаратное обеспечение SDN-балансировщика}





  SDN-балансировщик нагрузки (далее балансировщик) включает два 


основных компонента: передающее устройство (OpenFlow-коммутатор) 


и~SDN-конт\-рол\-лер. Далее отдельно рассмотрим каждый из них.





\subsection{Передающее устройство}





  Задачей передающего устройства является коммутация приходящих  


Ethernet-фрей\-мов на основании таблицы потоков, полученной от контроллера. 


Самостоятельными интеллектуальными функциями обработки трафика данное 


устройство не обладает. Взаимодействие с~контроллером осуществляется по 


протоколу OpenFlow или схожим. Программно-аппаратная реализация 


передающего устройства может различаться в~зависимости от требуемой 


пропускной способности.


  


  Если \textbf{объем балансируемого трафика не превышает~10~ГБ\!/\!с  


Full-Duplex}, то в~качестве передающего устройства можно использовать 


сервер общего назначения с~двумя портами~10G. В~экспериментах, описанных 


в~данной статье, использовалось передающее устройство со следующими 


основными характеристиками:


  \begin{itemize}


\item процессор: 1~процессор Intel Xeon E5-2620~v3 (2,40~ГГц, 6~ядер);


\item оперативная память: 8~ГБ;


\item сетевые адаптеры 10G (полезный трафик): 2~адаптера Intel 82599ES;


\item сетевые адаптеры 1G (OpenFlow, управление): 2~адаптера Intel  


I350-AM2.


\end{itemize}


  


  Обязательное требование к~сетевым адаптерам 10G~--- поддержка 


технологии Intel DPDK (Data Plane Development Kit,


далее DPDK). Тип и~объем жесткого диска критического значения не 


имеют, так как в~ходе работы OpenFlow-ком\-му\-та\-то\-ра регулярных 


операций чтения-записи не происходит. Следует добавить, что приведенная 


выше аппаратная конфигурация не является минимальной для  


SDN-ба\-лан\-си\-ров\-ки трафика~10~ГБ\!/\!с Full-Duplex. В~задачи данной 


работы не входила оптимизация аппаратных ресурсов с~целью максимального 


снижения стоимости передающего устройства.


  


  Основное программное обеспечение экспериментального передающего 


устройства включало Debian Server~8.7.1, Open vSwitch~2.7.0 и~Intel 


DPDK~16.11.


  


  Во всех экспериментах, описанных в~данной статье, на передающем 


устройстве использовалась технология  DPDK для 


предоставления программному коммутатору Open vSwitch (далее OVS), 


функционирующему в~пользовательском окружении Linux, эксклюзивного 


доступа к~ресурсам. Благодаря этому была получена производительность 


коммутатора OVS 10~ГБ\!/\!с Full-Duplex для IP-па\-ке\-тов с~MTU (maximum transmission unit)


в~диапазоне~256--1500~Б. Для IP-па\-ке\-тов размером~64~Б максимальная 


производительность OVS составила~6990~МБ\!/\!с Half-Duplex. Нагрузочное 


тестирование проводилось при помощи инструмента Pktgen~3.1.2. В~ходе 


настройки OVS\!/\!DPDK для каждого порта~10G создавались четыре  


RX-оче\-ре\-ди. Каждая очередь ассоциировалась с~одним логическим ядром 


процессора. Для портов~10G использовался драйвер vfio-pci. Также для работы 


OVS\!/\!DPDK на этапе загрузчика Linux резервировались~2048 2M-стра\-ниц 


оперативной памяти.


  


  Если \textbf{объем балансируемого трафика превышает 10~ГБ\!/\!с  


Full-Duplex}, то в~качестве передающего устройства можно использовать 


аппаратный коммутатор с~требуемым числом портов~10\!/\!40\!/\!100G и~поддержкой 


протокола OpenFlow нужной версии.





\vspace*{-2pt}





\subsection{SDN-контроллер}





  Аппаратное обеспечение SDN-конт\-рол\-ле\-ра может значительно 


различаться в~зависимости от максимально возможной нагрузки на контроллер. 


Нагрузка на контроллер зависит от количества новых потоков в~секунду 


(OpenFlow-со\-об\-ще\-ния PACKET\_IN от передающих устройств), которые 


контроллер должен обработать, и~от сложности логики обработки каждого 


потока. Ключевой аппаратной характеристикой для работы контроллера 


является производительность процессора. В~экспериментах на сервере 


контроллера использовались 2~процессора Intel Xeon E5-2609~v3 (1,90~ГГц, 


6~ядер), что было избыточным для нагрузки, создаваемой на контроллере 


в~ходе тестирования. Для управ\-ля\-ющей связи с~передающими устройствами 


использовался порт~1G. Программное обеспечение контроллера включало: 


Debian Server~8.7.1 и~модифицированный контроллер Beacon.


  


  Открытый контроллер Beacon написан на языке Java и~обладает следующими 


преимуществами по сравнению со многими другими  


SDN-конт\-рол\-ле\-ра\-ми~\cite{7-sdn, 8-sdn}: многопоточность, высокая 


производительность, легкость разработки собственных приложений. Основные 


недостатки контроллера Beacon: старшая поддерживаемая версия протокола 


OpenFlow-1.0, нестабильность работы, в~настоящее время не развивается 


и~не поддерживается. Модификация исходного контроллера Beacon 


заключалась в~минимизации его функционала до самого необходимого для 


повышения стабильности работы, а~также в~разработке собственных 


приложений для балансировки нагрузки.





\vspace*{-4pt}


  


\section{Отказоустойчивость SDN-балансировщика нагрузки}





\vspace*{-2pt}





  Отказоустойчивость SDN-балансировщика нагрузки зависит от следующих 


факторов:


  \begin{itemize}


\item отказоустойчивость при обрыве канала связи, по которому происходит 


балансировка трафика;\\[-14pt]


\item отказоустойчивость контроллера;\\[-14pt]


\item отказоустойчивость передающих устройств.


\end{itemize}





  Далее рассмотрим варианты реализации первых двух позиций.


  


\subsection{Отказоустойчивость при~обрыве канала связи}





  Алгоритм обеспечения отказоустойчивости при обрыве канала связи, по 


которому происходит балансировка трафика, рассмотрен на примере схемы 


масштабируемой защищенной связи двух географически распределенных 


корпоративных се\-тей/ЦОД. На рис.~2 представлены физическая и~логическая 


топологии распределенной сети. Далее приводится описание алгоритма 


отказоустойчивости.





\begin{figure}[b] %fig2


 \vspace*{1pt}


 \begin{center}


 \mbox{%


 \epsfxsize=127mm 


 \epsfbox{chsd-2.eps}


 }


\end{center}


 \vspace*{-11pt}


\Caption{Физическая и~логическая топологии для алгоритма отказоустойчивости при обрыве 


канала связи: \textit{1}~--- балансируемый корпоративный трафик; \textit{2}~--- 


управляющий OpenFlow-трафик между контроллером и~передающим устройствами; 


\textit{3}~--- GRE-тон\-не\-ли между передающими устройствами для передачи  


LLDP-со\-об\-ще\-ний %; ПУ~--- передающее устройство


}


\end{figure}


  


  На рис.~2 показана схема соединения двух сегментов корпоративной сети 


через публичную сеть при помощи масштабируемого защищенного канала 


связи. Шифрование и~расшифрование корпоративного трафика выполняются 


при помощи шести L3-крип\-то\-шлю\-зов, установленных по три на границе 


каждого сегмента корпоративной сети. Балансировка открытого 


корпоративного трафика на крип\-то\-шлю\-зы осуществляется на передающих 


устройствах (ПУ~1 и~2), управляемых контроллером по протоколу OpenFlow. 


Контроллер располагается в~первом сегменте корпоративной сети 


и~связывается с~передающим устройством второго сегмента при помощи 


отдельного защищенного канала связи. Между пе\-ре\-да\-ющи\-ми устройствами 


созданы GRE-туннели в~количестве, равном числу каналов, по которым 


происходит балансировка. В примере используются три канала. Каждый 


туннель проходит через <<свою>> пару криптошлюзов. GRE-тун\-не\-ли 


необходимы для создания L2-сег\-мен\-тов между передающими устройствами 


для передачи сообщений LLDP (Link Layer Discovery Protocol). 


Каждое передающее устройство 


периодически (в~эксперименте~--- каждые~5~с) посылает во\linebreak\vspace*{-12pt}





\pagebreak





\





\vspace*{-12pt}





\begin{floatingfigure}{60mm} %fig3


 \vspace*{-3pt}


 \begin{center}


 \mbox{%


 \epsfxsize=55.63mm 


 \epsfbox{chsd-3.eps}


 }


\end{center}


 \vspace*{-13pt}


\Caption{Обрыв канала связи}


\end{floatingfigure}





\noindent


 все свои  


GRE-тун\-не\-ли LLDP-со\-об\-ще\-ния, в~которых помимо прочей служебной 


информации содержатся идентификатор устройства и~идентификатор порта, 


с~которого было отправлено сообщение. Все полученные от <<соседа>>  


LLDP-со\-об\-ще\-ния передающее устройство пересылает на контроллер. 


Таким образом, контроллер в~режиме реального времени <<знает>> состояние 


всех каналов связи между передающими устройствами. 





Предположим, на 


одном из каналов произошла авария (рис.~3).


    В этом случае по одному из GRE-туннелей прекращается передача  


LLDP-со\-об\-ще\-ний. В~результате контроллер не получает никакой новой 


информации об \mbox{LLDP-тра}\-фи\-ке, ассоциированном с~данным туннелем. После 


неполучения нескольких последовательных LLDP-со\-об\-ще\-ний 


(в~эксперименте~--- двух) контроллер на всех передающих устройствах 


удаляет из таблицы потоков все правила, отправляющие трафик через порты, 


терминирующие отказавший канал связи. Далее потоки трафика, которые ранее 


коммутировались через отказавший канал, становятся <<новыми>> для 


передающего устройства и~согласно стандартной логике SDN отправляются на 


контроллер для анализа. В~итоге контроллер перенаправляет их через другие 


работоспособные каналы. 





Описанное GRE-тун\-не\-ли\-ро\-ва\-ние выполняется 


только для LLDP-тра\-фи\-ка. Полезный трафик корпоративной сети передается 


балансировщиком на криптошлюзы без ка\-кой-ли\-бо дополнительной 


инкапсуляции.


  


\subsection{Отказоустойчивость контроллера}


  


  Контроллер является важнейшим и~критичным звеном SDN-ин\-фра\-струк\-ту\-ры, 


так как полностью определяет поведение передающих устройств. Контроллер 


использует L3-ка\-нал для управляющей связи с~передающими устройствами. 


Следовательно, для резервирования контроллера в~режиме Active-Passive имеет 


смысл использовать технологию с~общим для всех серверов контроллера 


виртуальным IP-ад\-ре\-сом, который будут <<знать>> все передающие 


устройства. На рис.~4 показана схема резервирования контроллера в~режиме 


Active-Passive. Далее приведено описание варианта алгоритма 


отказоустойчивости.





\begin{figure} %fig4


 \vspace*{1pt}


 \begin{center}


 \mbox{%


 \epsfxsize=102.372mm 


 \epsfbox{chsd-4.eps}


 }


\end{center}


 \vspace*{-11pt}


\Caption{Схема резервирования контроллера в~режиме Active-Passive: \textit{1}~--- 


корпоративный трафик; \textit{2}~--- OpenFlow; \textit{3}~--- VRRP Keepalived


%; ПУ~--- передающее устройство


}


\end{figure}





  В схеме, показанной на рис.~4, для резервирования контроллера в~режиме 


Active-Passive используется VRRP (Virtual Router Redundancy Protocol, 


Linux Keepalived). Активный 


и~пассивный контроллеры, передающее устройство и~L3-крип\-то\-шлюз, 


предназначенный для защиты управляющего OpenFlow-канала к~другим 


географическим локациям, подключены к~одной локальной сети 172.16.1.0\!/\!24. 


На интерфейсах активного и~пассивного контроллера настраивается 


VRRP-груп\-па, в~которой помимо прочего определяются: приоритеты контроллеров, 


общий виртуальный IP-адрес, таймеры. Активным (MASTER) становится 


контроллер с~наибольшим приоритетом (на рис.~4 $\mathrm{P}\hm = 255$). Все 


передающие устройства для связи с~контроллером используют виртуальный  


IP-адрес 172.16.1.254, а~отвечает на ARP (Address Resolution Protocol)


и~OpenFlow запросы от передающих 


устройств именно активный контроллер. Также активный контроллер 


регулярно (по умолчанию каждую секунду) посылает пассивному сообщение 


Keepalive. В~схеме на рис.~4 данное сообщение передается через те же 


интерфейсы, что и~сообщения протокола OpenFlow. Это исключает ситуацию, 


когда управляющий OpenFlow-ка\-нал в~результате аварии недоступен, 


а~выделенный канал для сообщений Keepalive функционирует нормально, 


в~результате чего аварийного переключения на пассивный контроллер не 


происходит. При неполучении пассивным контроллером нескольких 


последовательных сообщений Keepalive (по умолчанию в~течение~3--4~с) 


последний становится активным, обновляет ARP-кэш передающих устройств 


и~начинает отвечать на OpenFlow-со\-об\-ще\-ния от них. Для реализации 


отказоустойчивости помимо протокола VRRP на каждом  


сер\-ве\-ре-конт\-рол\-ле\-ре используется специальный процесс (на схеме 


Watchdog), который регулярно отслеживает наличие в~памяти и~состояние 


процесса SDN-конт\-рол\-ле\-ра и~при необходимости перезапускает его.





\section{Алгоритм балансировки и~результаты экспериментов}


  


  В данном разделе кратко рассмотрен алгоритм балансировки нагрузки, 


использованный во всех экспериментах, а также представлены результаты 


балансировки на такие узлы защищенных сетей, как L3-крип\-то\-шлюз,  


TLS-шлюз, IDS, в~качестве которых использовались продукты и~новые 


разработки компании ИнфоТеКС (ViPNet Coordinator HW, ViPNet  


TLS-Gateway и~ViPNet IDS). Однако предложенные алгоритмы могут быть с~тем 


же успехом применены и~для оборудования других производителей.





\subsection{Алгоритм балансировки нагрузки}


  


  Алгоритм балансировки нагрузки был реализован на языке 


программирования Java в~виде модуля контроллера Beacon. Работа алгоритма 


начинается с~чтения конфигурационного файла, в~котором содержатся:  


IP-ад\-ре\-са устройств, на которые происходит балансировка; IP-ад\-рес 


доступа самого балансировщика; метрики направлений балансировки. Данная 


операция выполняется один раз при инициализации модуля. Направлением 


балансировки может быть исходящий порт балансировщика или IP-ад\-рес 


подключенного к~нему устройства, требующего балансировки. Далее 


формируется однострочный массив, в~котором каждое направление 


балансировки встречается столько раз, сколько указывает его метрика (по 


умолчанию~100). Чтобы распределение потоков трафика по направлениям 


балансировки было равномерным, среднее число новых потоков в~секунду 


должно превышать сумму метрик всех направлений балансировки. В~качестве 


идентификатора потока использовались четыре параметра: IP-ад\-рес 


отправителя, IP-ад\-рес получателя, транспортный порт отправителя, 


транспортный порт получателя.


  


  При инициализации передающего устройства контроллер загружает в~его 


таблицу потоков правило по умолчанию~--- все новые потоки отправлять на 


контроллер. Далее для каждого нового потока (при получении контроллером 


сообщения PACKET\_IN) контроллер случайным образом выбирает из массива 


направление балансировки и~инсталлирует новое правило в~таблицу потоков 


передающего устройства. Каждое правило имеет таймер неактивности (Idle 


Timeout, в~экспериментах~--- 60~с), благодаря которому правило 


автоматически удаляется из таблицы потоков при отсутствии трафика данного 


потока в~течение~60~с. Дополнительно для каждого потока могут выполняться 


операции изменения содержимого полей служебных заголовков 


фрей\-ма\!/\!па\-ке\-та\!/\!сег\-мен\-та, отправка потока в~туннель GRE\!/\!VXLAN и~др. 


В~рамках алгоритма также был реализован механизм ARP-Responder, 


позволяющий генерировать ответы на ARP-за\-про\-сы, приходящие на 


передающее устройство.


  


\subsection{Балансировка нагрузки на~L3-криптошлюзы}


  


  Рассмотрим результаты тестирования балансировки нагрузки на  


L3-крип\-то\-шлю\-зы. На рис.~5 представлена схема стенда.





\begin{figure} %fig5


 \vspace*{1pt}


 \begin{center}


 \mbox{%


 \epsfxsize=108.933mm 


 \epsfbox{chsd-5.eps}


 }


\end{center}


 \vspace*{-11pt}


\Caption{Схема стенда для тестирования балансировки нагрузки на L3-крип\-то\-шлюзы}


\end{figure}


  


  При проведении нагрузочного тестирования для генерации тестового 


трафика использовался инструмент Pktgen~3.1.2. Передающие устройства 


(ПУ~1 и~2) подключены к~Pktgen-сер\-ве\-ру при помощи 


оптических~10G~сетевых адаптеров Intel 82599ES, на которых активирована 


технология DPDK.


  


  Обозначения HW1000 и~HW5000 соответствуют L3-крип\-то\-шлю\-зам 


ViPNet Coordinator HW1000 и~ViPNet Coordinator HW5000 соответственно. 


Пропускная способность шифрованного канала HW1000, использованного 


в~экспериментах, составляет~900--1000~МБ\!/\!с, HW5000~--- 6--6,5~ГБ\!/\!с. 


Используемый алгоритм шифрования~--- ГОСТ28147-89 в~режиме CBC


(Cipher Block Chaining). 


Поскольку исследование производительности L3-крип\-то\-шлю\-зов и~прочих 


использованных в~работе устройств защищенных сетей не входило в~задачи 


данной статьи, подробная информация об их про\-грам\-мно-ап\-па\-рат\-ном 


обеспечении и~функционале здесь не приводится и~частично может быть 


найдена в~\cite{9-sdn}. Тестовый UDP-тра\-фик передавался в~режиме  


Half-Duplex (порт-от\-пра\-ви\-тель Pktgen-сер\-ве\-ра~--- Port~0,  


порт-по\-лу\-ча\-тель~--- Port~1) и~включал~2000~потоков в~секунду, 


полученных путем перебора транспортного порта отправителя 


в~диапазоне~10001--12000. Данного количества потоков было достаточно для 


тестирования равномерности балансировки нагрузки. Задача исследования 


максимальной производительности контроллера в~данной работе не стояла 


и~уже была решена в~\cite{8-sdn, 9-sdn}. Для алгоритма балансировки были 


выбраны следующие метрики направлений балансировки: HW1000~--- 


метрика~10, HW5000~--- метрика~60.


  


  В таблице представлены результаты тестирования балансировки нагрузки на 


L3-крип\-то\-шлю\-зы. Под пропускной способностью понимается объем 


тестового трафика в~секунду, приходящего на порт~1 Pktgen-сер\-ве\-ра после 


отправки с~порта~0, балансировки на передающем устройстве ПУ~1, 


шифрования и~рас\-шиф\-ро\-ва\-ния на крип\-то\-шлю\-зах HW1000\!/\!5000 и~сборки на 


передающем устройстве ПУ~2. В~ходе всех тестов нагрузка на  


L3-крип\-то\-шлю\-зы была максимально возможной. Загрузка процессоров 


криптошлюзов составляла~98\%--100\%.


  


  \begin{table}\small


  \begin{center}


  \tabcolsep=1.7pt


  \begin{tabular}{|l|c|c|c|}


  \multicolumn{4}{c}{Результаты тестирования балансировки нагрузки на L3-криптошлюзы}\\


  \multicolumn{4}{c}{\ }\\[-6pt]


  \hline


\multicolumn{1}{|c|}{L3-криптошлюзы, на которые}&


 \multicolumn{3}{c|}{Пропускная способность, 


МБ\!/\!с}\\


\cline{2-4}


\multicolumn{1}{|c|}{\tabcolsep=0pt


\begin{tabular}{c}направлена балансировка нагрузки\\


 (одна сторона защищенного канала)\end{tabular}}&\tabcolsep=0pt


\begin{tabular}{c}UDP,\\


$\mathrm{MTU}= 1432$~Б\end{tabular}&


\tabcolsep=0pt


\begin{tabular}{c}UDP,\\ $\mathrm{MTU}=256$~Б\end{tabular}&


\tabcolsep=0pt


\begin{tabular}{c}UDP,\\ $\mathrm{MTU}=64$~Б\end{tabular}\\


\hline


1 HW1000&\hphantom{9}933&\hphantom{9}461&180\\


1 HW5000&6241&1016&288\\


3 HW1000 (балансировка)&2806&1398&552\\


3 HW1000\;+\;1 HW5000 (балансировка)&\textbf{9041}&\textbf{2419}&\textbf{858}\\


\hline


\end{tabular}


\end{center}


\end{table}


  


  В первой колонке таблицы перечислены L3-крип\-то\-шлю\-зы, 


подключенные к~одному передающему устройству. Передающее устройство на 


противоположном конце защищенного канала балансирует\!/\!принимает трафик 


на\!/\!от аналогичного набора L3-крип\-то\-шлю\-зов (см.\ рис.~5). Как видно из 


таблицы, пропускная способность шифрованного канала при балансировке 


трафика на три криптошлюза HW1000 и~один криптошлюз HW5000 


практически равна сумме пропускных способностей отдельных криптошлюзов. 


Небольшие колебания значений вызваны спецификой работы инструмента 


Pktgen.





\subsection{Балансировка нагрузки на~TLS-шлюзы}





  На рис.~6 представлена схема стенда для тестирования балансировки 


нагрузки на TLS-шлюзы.





\begin{figure}[b] %fig6


 \vspace*{1pt}


 \begin{center}


 \mbox{%


 \epsfxsize=128.742mm 


 \epsfbox{chsd-6.eps}


 }


\end{center}


 \vspace*{-11pt}


\Caption{Схема стенда для тестирования балансировки нагрузки на TLS-шлюзы}


\end{figure}





  Схема на рис.~6 состоит из следующих основных компонентов:


  \begin{itemize}


\item TLS-Meter~--- генератор тестового трафика, представляющего собой 


HTTPS-за\-про\-сы, идущие на виртуальный адрес передающего устройства 


172.16.1.250. Данный инструмент разработан в~компании ИнфоТеКС, 


в~описываемых экспериментах использует алгоритм ГОСТ~Р~34.10-2012;


\item SDN-балансировщик, включающий передающее устройство 


и~контроллер;


\item TLS-шлюзы (TLS-Gw~1--4), на которые происходит балансировка 


HTTPS-за\-про\-сов. Все четыре шлюза были представлены одной тестовой 


моделью, работающей в~режиме HTTPS-Proxy с~использованием алгоритма 


шифрования ГОСТ~Р~34.10-2012;


\item Web-сервер~--- веб-сервер, которому TLS-шлю\-зы адресуют 


расшифрованные HTTP-за\-просы.


  \end{itemize}


  


  \begin{figure}[b] %fig7


 \vspace*{1pt}


 \begin{center}


 \mbox{%


 \epsfxsize=125.591mm 


 \epsfbox{chsd-7.eps}


 }


\end{center}


 \vspace*{-11pt}


\Caption{График зависимости количества установленных HTTPS-со\-еди\-не\-ний от 


времени теста для каждого TLS-шлю\-за (\textit{1}--\textit{4}~--- TLS-Gw~1--TLS-Gw~4)


и~суммарная балансируемая нагрузка~(\textit{5})}


\end{figure}


  Генерация тестового трафика выполнялась с~240~IP-ад\-ре\-сов. 


Длительность теста~--- 3~мин, версия TLS-1.2. В~ходе тестов измерялось 


количество установленных HTTPS-со\-еди\-не\-ний в~секунду в~зависимости от 


числа TLS-шлю\-зов. Полученные результаты не являются предельными 


характеристиками, так как нагрузочное тестирование TLS-шлюза в~задачи 


данной работы не входило. Замеры нагрузки производились на  


L2-ком\-му\-та\-то\-ре, соединяющем передающее устройство и~TLS-шлю\-зы, 


путем SPAN (Switched Port ANalyzer) зер\-ка\-ли\-ро\-ва\-ния всего проходящего через коммутатор 


трафика на выделенный персональный компьютер, сбора трафика при помощи tcpdump и~анализа при 


помощи Wireshark. На рис.~7 показаны графики зависимости количества 


установленных HTTPS-со\-еди\-не\-ний от времени теста для каждого  


TLS-шлю\-за, а также суммарная нагрузка.





\pagebreak





\





\vspace*{-12pt}











\begin{floatingfigure}{68mm} %fig8


 \vspace*{-1pt}


 \begin{center}


 \mbox{%


 \epsfxsize=62.961mm 


 \epsfbox{chsd-8.eps}


 }


\end{center}


 \vspace*{-11pt}


\Caption{График зависимости количества установленных HTTPS-со\-еди\-не\-ний в~секунду 


от числа TLS-шлю\-зов}


\end{floatingfigure}





  


  Как видно на рис.~7, производительность каждого отдельного TLS-шлю\-за 


в~ходе всего теста со\-став\-ля\-ла приблизительно~235~соединений в~секунду. Для 


всех четырех\linebreak TLS-шлю\-зов наблюдались равномерность и~одинаковость 


загрузки, что говорит о корректной балансировке HTTPS-тра\-фи\-ка. 


Суммар\-ная\linebreak производительность системы со\-ста\-ви\-ла четырехкратное увеличение 


производительности одного TLS-шлю\-за. 





На рис.~8 показан график 


зависимости количества установленных HTTPS-со\-еди\-не\-ний в~секунду от 


чис\-ла TLS-шлю\-зов, нагрузка на которые проходила балансировку.


      График, представленный на рис.~8, имеет почти линейный характер, что 


говорит о~том, что количество на\-прав\-ле\-ний балансировки практически не 


влияет на производительность балансировщика. Поскольку в~тестах 


использовались TLS-шлю\-зы одинаковой производительности, метрики 


на\-прав\-ле\-ний балансировки имели одинаковые значения~--- 100. В~данном 


эксперименте для реализации от\-ка\-зо\-устой\-чи\-вости каналов связи, по которым 


происходила балансировка, использовался сле\-ду\-ющий алгоритм. 


Каждые~2~с контроллер генерирует с~пе\-ре\-да\-юще\-го устройства  


ARP-за\-про\-сы на IP-ад\-ре\-са всех TLS-шлю\-зов  


(172.16.1.251--172.16.1.254). В~случае неполучения двух последовательных 


ARP-от\-ве\-тов от одного TLS-шлю\-за данное направление балансировки 


считается нефункциональным и~исключается из алгоритма балансировки. При 


этом отправка ARP-запросов на IP-ад\-рес <<выбывшего>> TLS-шлю\-за 


продолжается, и~в~случае получения ARP-от\-ве\-та данное направление 


балансировки снова начинает использоваться для передачи полезного трафика.


  


\subsection{Балансировка нагрузки на~системы обнаружения вторжений}





  На рис.~9 представлена схема стенда для тестирования балансировки 


нагрузки на масштабируемую систему обнаружения вторжений (IDS).





\begin{figure} %fig9


 \vspace*{1pt}


 \begin{center}


 \mbox{%


 \epsfxsize=104.626mm 


 \epsfbox{chsd-9.eps}


 }


\end{center}


 \vspace*{-11pt}


\Caption{Схема стенда для тестирования балансировки нагрузки на масштабируемую 


систему обнаружения вторжений}


\end{figure}


  


  В схеме на рис.~9 полезный трафик зеркалируется на L2-ком\-му\-та\-то\-ре 


и~отправляется на передающее устройство SDN-балансировщика, где 


равномерно балансируется на три системы обнаружения вторжений, 


максимальная производительность каждой из которых составляет~6~ГБ\!/\!с 


(использовалась модель ViPNet IDS2000). Особенностью алгоритма 


балансировки для данного случая является то, что при обработке каждого 


нового потока контроллер инсталлирует в~таблицу потоков передающего 


устройства два правила коммутации: прямое и~обратное. В~прямом правиле  


IP-ад\-ре\-са и~транспортные порты отправителя и~получателя соответствуют 


полученным из первого пакета нового потока. В~обратном правиле 


идентификаторы отправителя и~получателя меняются местами. Оба правила 


используют одно и~то же направление балансировки. Таким образом 


гарантируется, что весь трафик одной TCP/UDP-сес\-сии будет анализироваться 


одной IDS.


  


  Кратко рассмотрим результаты эксперимента. В~ходе теста загрузка сетевых 


интерфейсов всех IDS была практически одинаковой и~синхронно менялась при 


изменении объема подаваемого на передающее устройство трафика. Анализ 


трафика, приходящего на каждую IDS после балансировки, показал отсутствие 


<<разорванных>> TCP-сес\-сий, что говорит о том, что контроллер корректно 


распределяет потоки трафика по направлениям балансировки.





\section{Заключение}





  Технология SDN благодаря централизованному управлению сетью 


и~программируемости хорошо подходит для решения задач балансировки 


трафика\linebreak


 при необходимости согласованной работы балансировщиков, 


от\-ка\-зо\-устой\-чи\-вости и~нестандартных алгоритмов балансировки. Open vSwitch 


при использовании технологии Intel DPDK успешно коммутирует не 


менее~10~ГБ\!/\!с Full-Duplex-тра\-фи\-ка. Предложены и~частично 


протестированы механизмы от\-ка\-зо\-устой\-чи\-вости каналов связи, по которым 


происходит балансировка (режим Active-Active), и~SDN-конт\-рол\-лера (режим 


Active-Passive). Успешно проведены экспериментальные исследования 


балансировки нагрузки на такие устройства защищенных сетей, как L3-крип\-то\-шлюз, 


TLS-шлюз, IDS. Во всех экспериментах был использован 


унифицированный алгоритм балансировки, использующий весовые 


коэффициенты каналов.


  


{\small\frenchspacing


{%\baselineskip=10.8pt


\addcontentsline{toc}{section}{Литература}


\begin{thebibliography}{99}





\bibitem{3-sdn} %1


\Au{Feamster N., Balakrishnan~H.} Detecting BGP configuration faults with static 


analysis~/\!\!/ 2nd Conference on Symposium on Networked Systems


Design and Implementation Proceedings, 2005. Vol.~2. 


P.~43--56.


\bibitem{4-sdn} %2


\Au{Sherry J., Ratnasamy~S.} A~survey of enterprise middlebox deployments.~--- 


Berkeley, CA, USA: Univ. California, Electr. Eng. Comput. Sci. Dept., 


2012. Technical Report  UCB/EECS-2012-24.


\bibitem{1-sdn} %3


\Au{Kim H., Feamster~N.} Improving network management with software defined 


networking~/\!\!/ IEEE Commun. Mag., 2013. Vol.~51. No.\,2. P.~114--119.


\bibitem{2-sdn} %4


\Au{Kreutz D., Ramos~F.\,M.\,V., Verissimo~P., Rothenberg~C.\,E., Azodolmolky~S., Uhlig~S.} 


Software-defined networking: A~comprehensive survey~/\!\!/ P.~IEEE, 2015. 


Vol.~103. No.\,1. P.~14--76.





\bibitem{5-sdn} %5


\Au{Назаров М.\,А.} Определение, основные понятия и~архитектуры 


про\-грам\-мно-кон\-фи\-гу\-ри\-ру\-емых сетей~--- SDN (Software Defined Networking)~/\!\!/ Информатизация 


и~связь, 2015. №\,4. С.~82--87.


\bibitem{6-sdn}


\Au{Greene K.} TR10: Software-defined networking~/\!\!/ 10 Breakthrough Technologies: MIT 


Technology Review, 2009. {\sf  


http://www2.technologyreview.com/article/412194/ tr10-software-defined-networking}. 


\bibitem{7-sdn}


\Au{Erickson D.} The Beacon OpenFlow controller~/\!\!/ 2nd ACM SIGCOMM Workshop on 


Hot Topics in Software Defined Networking Proceedings, 2013. P.~13--18.


\bibitem{8-sdn}


\Au{Shalimov A., Zuikov D., Zimarina~D., Pashkov~V., Smeliansky~R.} Advanced study of 


SDN/OpenFlow controllers~/\!\!/ 9th Central \& Eastern European Software Engineering 


Conference in Russia Proceedings, 2013. Article No.\,1.


{\sf 


https://www.researchgate. net/publication/262155093\_Advanced\_study\_of\_SDNOpenFlow\_controllers/overview}. 


\bibitem{9-sdn}


Infotecs: Продукты. {\sf http://infotecs.ru/product/all/?line=vipnet-network-security}.





        \end{thebibliography}


} }








\vspace*{-3pt}





\hfill{\small\textit{Поступила в~редакцию 26.11.17}}














%\vspace*{8pt}





%\hrule





%\vspace*{2pt}





\newpage





\vspace*{-28pt}





%\hrule





%\vspace*{9pt}











\def\tit{SDN LOAD BALANCING FOR~SECURE NETWORKS}





\def\titkol{SDN load balancing for~secure networks}








\def\aut{O.\,Yu.~Guzev$^1$ and I.\,V.~Chizhov$^{2,3}$}


\def\autkol{O.\,Yu.~Guzev and I.\,V.~Chizhov}











\titel{\tit}{\aut}{\autkol}{\titkol}





\vspace*{-9pt}








\noindent


$^1$Research and Development Center, JSC ``InfoTeCS,''  1\!/\!23, b.~1 Staryy\linebreak


$\hphantom{^1}$Petrovsko-Razumovskiy Pr., Moscow 127287, Russian Federation








\noindent


$^2$Faculty of Computational Mathematics and Cybernetics, M.\,V.~Lomonosov\linebreak


$\hphantom{^1}$Moscow State University, 2nd Education Building, Faculty CMC, GSP-1,\linebreak


$\hphantom{^1}$Leninskie  Gory, Moscow 119991, Russian Federation





\noindent


$^3$Institute of informatics Problems, Federal Research Center ``Computer Science\linebreak


$\hphantom{^1}$and Control'' of the Russian 


Academy of Sciences, 44-2~Vavilov Str., Moscow\linebreak


$\hphantom{^1}$119333, Russian Federation








\def\leftfootline{\small{\textbf{\thepage}


\hfill Sistemy i Sredstva Informatiki~---


Systems and Means of Informatics 2018 vol~28 no\,1}


}%


 \def\rightfootline{\small{Sistemy i Sredstva Informatiki~---


 Systems and Means of Informatics   2018   vol~28   no\,1


\hfill \textbf{\thepage}}}











\Abste{The software-defined networking (SDN) technology in comparison with traditional IP networks allows programming the 


network's behavior using a centralized controller. In this case, forwarding devices deal only with 


forwarding frames based on flow tables loaded into them by the controller. Flow tables are built on 


the controller during the processing of information about traffic flows arriving at forwarding 


devices. The above properties of the technology were used to create the SDN load balancer for 


devices of secure networks. The article discusses the architecture and software of the balancer. 


Descriptions of schemes and results of experiments on load balancing for 


such devices as L3-VPN (Level~3 Virtual Private Network)


gateway,  


TLS (Transport Layer Security) gateway, and IDS (Intrusion Detection System)


are given.}





\KWE{software-defined networking (SDN); controller; OpenFlow; VPN gateway; TLS; intrusion 


detection system; IDS; load balancing; DPDK; Open vSwitch; Beacon}








\DOI{10.14357\!/\!086965271801010}





%\vspace*{-24pt}





%\Ack


%\noindent





 








\renewcommand{\bibname}{\protect\rmfamily References}


%\renewcommand{\bibname}{\large\protect\rm References}





{\small\frenchspacing


{%\baselineskip=10.8pt


\addcontentsline{toc}{section}{References}


\begin{thebibliography}{99}





 \bibitem{3-sdn-1} %1


  \Aue{Feamster, N., and H. Balakrishnan.} 2005. Detecting BGP configuration faults with static 


analysis. \textit{2nd Conference on Symposium on Networked Systems 


Design and Implementation Proceedings}. 2:43--56.


  \bibitem{4-sdn-1} %2


  \Aue{Sherry, J., and S. Ratnasamy.} 2012. A~survey of enterprise middlebox deployments: 


 Berkeley, CA: Univ. California, Electr. Eng. Comput. Sci. 


Dept. Technical Report UCB/EECS-2012-24.


  \bibitem{1-sdn-1} %3


  \Aue{Kim, H., and N. Feamster.} 2013. Improving network management with software defined 


networking. \textit{IEEE Commun. Mag.} 51(2):114--119.


  \bibitem{2-sdn-1} %2


  \Aue{Kreutz, D., F.\,M.\,V.~Ramos, P.~Verissimo, C.\,E.~Rothenberg, S.~Azodolmolky, and 


S.~Uhlig.} 2015. Software-defined networking: A~comprehensive survey. 


\textit{P.~IEEE} 103(1):14--76.


 


  \bibitem{5-sdn-1}


  \Aue{Nazarov, M.\,A.} 2015. Opredelenie, osnovnye ponyatiya i~arkhitektury  


programmno-konfiguriruemykh setey~--- SDN (Software Defined Networking) [Determination, the 


basic concepts, and architecture of the program defined networks].  \textit{Informatizatsiya i~svyaz'} 


[Informatization and Communication] 4:82--87.


  \bibitem{6-sdn-1}


  \Aue{Greene, K.} 2009. TR10: Software-defined networking. \textit{10~Breakthrough 


Technologies: MIT Technology Review}. Available at: {\sf  


http://www2.technologyreview.com/ article/412194/tr10-software-defined-networking}  (accessed 


February~24, 2018).


  \bibitem{7-sdn-1}


  \Aue{Erickson, D.} 2013. The Beacon OpenFlow controller. \textit{2nd ACM SIGCOMM 


Workshop on Hot Topics in Software Defined Networking Proceedings}. 13--18.


  \bibitem{8-sdn-1}


  \Aue{Shalimov, A., D.~Zuikov, D.~Zimarina, V.~Pashkov, and R.~Smeliansky}. 2013. 


Advanced study of SDN/OpenFlow controllers. \textit{9th Central \& Eastern European Software 


Engineering Conference in Russia Proceedings}. Article No.\,1. Available at: {\sf 


https://www.researchgate.net/publication/262155093\_Advanced\_ study\_of\_SDNOpenFlow\_controllers/overview} 


(accessed February~24, 2018).


  \bibitem{9-sdn-1}


  Infotecs. Produkty [Products]. Available at:


  {\sf http://infotecs.ru/product/all/?line= vipnet-network-security}  


(accessed February~24, 2018).


\end{thebibliography}


} }








\vspace*{-3pt}





\hfill{\small\textit{Received November 26, 2017}}   


   


   \Contr


   


\noindent


\textbf{Guzev Oleg Yu.} (b.\ 1980)~--- Candidate of Science (PhD) in technology, 


researcher, Research and Development Center, JSC ``InfoTeCS,'' 1/23, b.~1, 


Staryy Petrovsko-Razumovskiy Pr., Moscow 127287, Russian Federation; 


\mbox{oleg.guzev@infotecs.ru}





\vspace*{3pt}





\noindent


\textbf{Chizhov Ivan V.} (b.\ 1984)~--- Candidate of Science (PhD) in physics 


and mathematics, associate professor, Faculty of Computational Mathematics and 


Cybernetics, M.\,V.~Lomonosov Moscow State University, 2nd Education Building, 


Faculty CMC, GSP-1, Leninskie Gory, Moscow 119991, Russian Federation; senior 


scientist, Institute of Informatics Problems, Federal Research Center 


``Computer Science and Control'' of the Russian Academy of Sciences,  


44-2~Vavilov Str., Moscow 119333, Russian Federation; \mbox{ichizhov@cs.msu.ru}





\label{end\stat}








\renewcommand{\bibname}{\protect\rm Литература}


